\subsection{Irreducible elements}

DEFINITION 6. --- An element x of an ordered group G is said to be irreducible if it is a minimal element in the set of strictly positive elements of G.

Let x be an irreducible element of the ordered group G; if y is a positive element of G then the element \(\inf(x, y)\) , if it exists, can only be equal to x or to 0. Thus in a lattice ordered group G, every positive y is either greater than or coprime to the irreducible element x; in particular, two distinct irreducible elements are coprime.

\begin{enumerate}
\def\labelenumi{(\Roman{enumi})}
\setcounter{enumi}{503}
\tightlist
\item
  An element \(p\) of \(\mathbf{A}\) is called irreducible if the ideal \((p)\) is an irreducible element of the ordered group \(\mathcal{P}^*\) ; this says that \(p\) is neither zero nor invertible, and that any element of \(\mathbf{A}\) which divides \(p\) is associate either to \(p\) or to 1. If \(\mathcal{P}^*\) is lattice ordered, then every \(a \in \mathbf{A}\) is either coprime to \(p\) or a multiple of \(p\) .
\end{enumerate}

Examples (DIV). --- 1) An integer \(p > 0\) is irreducible in \(\mathbf{Z}\) if and only if it is prime (I, p.~50).

\begin{enumerate}
\def\labelenumi{\arabic{enumi})}
\setcounter{enumi}{1}
\tightlist
\item
  A polynomial in one indeterminate over a field K is irreducible if and only if it is irreducible in the usual sense (IV, p.~13).
\end{enumerate}

PROPOSITION 14. --- For an element \(x > 0\) of an ordered group \(G\) to be irreducible it is sufficient that it satisfy the following property:

\begin{enumerate}
\def\labelenumi{(\Alph{enumi})}
\setcounter{enumi}{15}
\tightlist
\item
  The relations \(x \leqslant y + z\) , \(y \geqslant 0\) , \(z \geqslant 0\) imply \(x \leqslant y\) or \(x \leqslant z\) .
\end{enumerate}

This condition is necessary when G is lattice ordered.

If G is lattice ordered and x is irreducible, we have just seen that y is either greater than x or coprime to x; in the latter case Cor. 2 of VI, p.15 shows that z is greater than x. Conversely, suppose the condition is satisfied: if \(0 \leqslant y \leqslant x\) then it follows, by putting \(x = y + z\) ( \(z \geqslant 0\) ), that either \(x \leqslant y\) or \(x \leqslant z\) ; in the first case we have x = y; in the second we have \(x \leqslant x - y\) , so \(y \leqslant 0\) and so y = 0; this shows that x is indeed irreducible.

PROPOSITION 14 (DIV). --- For a nonzero element p of A to be irreducible it is sufficient that it not be a unit, and that it cannot divide a product of two elements of A without dividing one or other of them. This condition is necessary if \(P^{*}\) is lattice ordered.

Remark. --- Proposition 14 (DIV) can also be expressed as follows: if p is a non zero element of A such that the ideal \((p)\) is prime (I, p.~117, Def. 3) then p is irreducible; conversely, if \(P^{*}\) is lattice ordered and p is irreducible then the ideal \((p)\) is prime.

PROPOSITION 15. --- In an ordered group G, let \((p_{\iota})_{\iota \in I}\) be a family of pairwise distinct positive elements of G satisfying condition \((P)\) (and hence irreducible). Then the map

\[
(n _ {i}) _ {i \in I} \mapsto \sum_ {i \in I} n _ {i} p _ {i}
\]

is an isomorphism of the ordered group \(\mathbf{Z}^{(I)}\) , the direct sum of the ordered groups Z (VI, p.~7) onto the ordered subgroup of G generated by the \(p_{i}\) .

It is enough to show that the relation \(\sum_{\iota \in I} n_{\iota} p_{\iota} \geqslant 0\) is equivalent to \(n_{\iota} \geqslant 0\) for all \(\iota\) ,

for in particular the relation \(\sum_{\iota \in I} n_{\iota} p_{\iota} = 0\) will imply \(n_{\iota} = 0\) for all \(\iota\) , hence this will

show that the family \((p_{\iota})\) is linearly independent. Now let \(I'\) (resp. \(I''\) ) be the finite subset of \(I\) consisting of those \(\iota\) such that \(n_{\iota} > 0\) (resp. \(n_{\iota} < 0\) ); we have

\[
\sum_ {i \in I ^ {\prime}} n _ {i} p _ {i} \geqslant \sum_ {i \in I ^ {\prime \prime}} (- n _ {i}) p _ {i}.
\]

In particular, for \(\lambda \in I''\) , this implies that \(p_{\lambda} \leqslant \sum_{i \in I'} n_i p_i\) , and it follows by induction

from property (P) that we must have \(p_{\lambda} \leqslant p_{\iota}\) for some \(\iota \in I'\) ; since \(p_{\iota}\) is irreducible, this would imply \(p_{\lambda} = p_{\iota}\) , which is absurd. Hence \(I''\) is empty, which proves the proposition.

THEOREM 2. --- Let G be a filtered group. Then the following properties are equivalent :

\begin{enumerate}
\def\labelenumi{\alph{enumi})}
\item
  G is isomorphic to an ordered group of the form \(\mathbf{Z}^{(I)}\) .
\item
  G is lattice ordered and satisfies the following condition:
\end{enumerate}

(MIN) Every nonempty set of positive elements of \(\mathbf{G}\) has a minimal element.

\begin{enumerate}
\def\labelenumi{\alph{enumi})}
\setcounter{enumi}{2}
\item
  G satisfies condition (MIN) and every irreducible element of G has property (P).
\item
  G is generated by its irreducible elements, and every irreducible element of G has property (P).
\end{enumerate}

Let us show first that \(a\) implies \(b\) . The group \(\mathbf{Z}^{(I)}\) is lattice ordered, as the direct sum of totally ordered groups. On the other hand let E be a nonempty set of positive elements of \(\mathbf{Z}^{(I)}\) and let \(x = \sum_{i} n_i e_i\) be an element of E (where

\((e_{i})\) denotes the natural basis of \(\mathbf{Z}^{(I)}\) ; there are a finite number \(\prod_{i}(n_{i} + 1)\) of

elements \(y\) of \(\mathbf{Z}^{(I)}\) such that \(0 \leqslant y \leqslant x\) , so the set \(F\) of elements of \(E\) which are less than or equal to \(x\) is a fortiori finite; since it is nonempty, it contains a minimal element (Set Theory, III, p.~170, Cor. 2), which is clearly a minimal element of \(E\) .

It is clear that \(b)\) implies \(c)\) , by Prop. 14. Let us show that \(c)\) implies \(d)\) . Since \(G\) is filtered, it is enough (VI, p.~4, Prop. 4) to check that the set \(F\) of positive elements of \(G\) which are sums of irreducible elements is equal to \(G_{+} - \{0\}\) . If this were not true, it would follow from (MIN) that the complement of \(F\) in \(G_{+} - \{0\}\) would have a minimal element \(a\) ; by definition \(a\) is not irreducible, so is the sum of two strictly positive elements \(x\) and \(y\) ; since \(x < a\) and \(y < a\) , these elements belong to \(F\) , and so are sums of irreducible elements, and it follows that so is \(a\) , which is a contradiction. Finally, \(d)\) implies \(a)\) by Prop. 15.

We will apply Th. 2 to the theory of divisibility in principal ideal domains (VII, p.~4) and in unique factorisation domains (AC, VII, § 3), as well as to the study of ideals in a Dedekind ring (AC, VII, § 2).

